\done\todo{Develop an abstract: The abstract summarizes the information contained in the rest of this document. An abstract includes:(a)the research question or engineering problem, (b) procedures used, (c) data, (d) interpretation, and (e) conclusions. It also may include any possible research applications. It should be limited to these essential elements. Since this is a proposal, the conclusion part should be a potential conclusion.}
 \\

\done\todo{Please address comments left on Overleaf.}
Software vulnerability detection is a crucial step in software development. High quality labeled vulnerability code is crucial in training Machine Learning (ML)/Artificial Intelligence (AI)-based accurate vulnerability detection models. However, the high cost and lack of human expertise makes large scale annotation a difficult task. This project investigates whether a small set of human annotation can improve accuracy in AI/ML vulnerability detection models. We leverage existing ML/AI models using uncertainty quantification and LLM filtering, to filter commit patches into a small enough subset(50-100), feasible for human annotation. Data for training will be obtained from open source machine learning projects for their high relevancy. Vulnerability detection models will be retrained with this small set of human annotated data in an attempt to display that a limited number of high quality annotation can yield statistically significant predictive performance improvements.